\subsection{Fixed determinant data with degenerating ellipticity}
Here $v_-(h)$ and $v_+(h)$ denote the infimum and supremum of
$\int_X(h+\ddc w)^n$ over bounded $h$-psh potentials $w$.
\begin{theorem}\label{thm:fixed-determinant-closed-degeneration}
On $X=(\C/(\mathbb Z+i\mathbb Z))^n$, $n\geq2$, there
are positive metrics $\omega_\varepsilon$ in a fixed K\"ahler
class, $0<\varepsilon\leq1$, with fixed complete current and
determinant-measure families and fixed endpoint \eqref{eq:constant-mass-determinant-profile},
while their least relative eigenvalue equals $\varepsilon$.
They have uniformly bounded smooth coefficients and global
potentials, and
\[
 v_-(\omega_\varepsilon)=v_+(\omega_\varepsilon)=n!>0.
\]
The manifold has a global holomorphic volume form and $c_1=0$.
Every displayed metric is already balanced and pluriclosed.
\end{theorem}
\begin{proof}
Write $z_j=x_j+iy_j$,
$\beta_j=(i/2)dz_j\wedge d\bar z_j=dx_j\wedge dy_j$,
$\beta=\sum_j\beta_j$, $V=n!$, and $\mathrm dV=\beta^n/V$.
Set
\[
 u_\varepsilon=\frac{1-\varepsilon}{4\pi}\cos(4\pi x_1),
 \qquad a_\varepsilon=1-(1-\varepsilon)\cos(4\pi x_1).
\]
For a function of $x_1$, $\ddc u=(u''/(4\pi))\beta_1$.
Hence
\[
 \omega_\varepsilon=\beta+\ddc u_\varepsilon
 =a_\varepsilon\beta_1+\sum_{j=2}^n\beta_j.
\]
These are globally periodic smooth positive closed forms with
$\varepsilon\leq a_\varepsilon\leq2-\varepsilon$ and
$\omega_\varepsilon\leq2\beta$. Their minimum relative eigenvalue
is $\varepsilon$. They converge smoothly to a semipositive form
losing rank on $\{x_1\in\tfrac12\mathbb Z\}$.
All fixed-order smooth norms of the coefficients and potentials are
uniformly bounded. Theorem~\ref{thm:constant-mass-entire-determinant-profile} proves equality of the
entire current data and the exact norm at every $\varepsilon$.

Since $\omega_\varepsilon$ is closed, \cref{thm:mass-invariance-characterization} gives mass $V$ for every bounded $\omega_\varepsilon$-plurisubharmonic potential. The explicit formula also gives $\int_X\omega_\varepsilon^n=V$, since $a_\varepsilon$ has mean one.
Thus both extrema over bounded
potentials equal $V$. The global form
$dz_1\wedge\cdots\wedge dz_n$ trivializes the canonical bundle;
closed positive metrics are balanced and pluriclosed.
\end{proof}

\subsection{The same obstruction for nonclosed pluriclosed metrics}
\begin{theorem}\label{thm:fixed-determinant-nonclosed-degeneration}
On the fixed torus above, fix $0<\tau<1/2$. There is a positive
nonclosed pluriclosed metric $h_\tau$ and a family
$\Omega_\varepsilon=h_\tau+\ddc u_\varepsilon>0$,
$0<\varepsilon\leq1/2$, with identical actual current sets and
intrinsic determinant-measure families. For any fixed smooth positive
probability volume their entire-family endpoint is
\[
 \sup_{T\in\mathscr G_{\Omega_\varepsilon}}
       \int_X|p|\,\Gamma_0(T)
       =V_\tau^{-1/n}\|p\|_{L^n(\mathrm dV)},\qquad
 V_\tau=n!(1-\tau^2/2).
\]
Both bounded-potential masses equal $V_\tau$. All fixed-order smooth
coefficient and potential norms are uniformly bounded, while
\[
 \min_X\lambda_{\min}(\Omega_\varepsilon\mid\beta)=\varepsilon,
 \qquad d\Omega_\varepsilon=dh_\tau\ne0.
\]
The flat $\beta$ is a fixed balanced metric on this same manifold
with $c_1=0$; each $\Omega_\varepsilon$ is pluriclosed.
\end{theorem}
\begin{proof}
With the torus coordinates above set $s=\sin(2\pi x_1)$ and
\[
 \xi=\frac i2(dz_1\wedge d\bar z_2+dz_2\wedge d\bar z_1)
      =dx_1\wedge dy_2+dx_2\wedge dy_1,
 \qquad h_\tau=\beta+\tau s\xi.
\]
The first relative Hermitian block is
$\left(\begin{smallmatrix}1&\tau s\\\tau s&1\end{smallmatrix}\right)$,
so positivity follows from $|\tau s|<1$. Moreover
\[
 dh_\tau=2\pi\tau\cos(2\pi x_1)
                   dx_1\wedge dx_2\wedge dy_1\ne0.
\]
Because $\ddc s$ is proportional to $dz_1\wedge d\bar z_1$,
$\ddc(s\xi)=0$. The first two-coordinate metric is thus pluriclosed,
and its square is top degree on that complex surface. Expanding with
the remaining closed curve factors proves
\[
 \ddc(h_\tau^j)=0\quad(1\leq j\leq n-1),\qquad
 \int_X h_\tau^n=n!\int_0^1(1-\tau^2\sin^2(2\pi x_1))dx_1
                   =V_\tau.
\]
Use the same $u_\varepsilon$ and $a_\varepsilon$ as in
Theorem~\ref{thm:fixed-determinant-closed-degeneration}, now with $\varepsilon\leq1/2$. Then
\[
 \Omega_\varepsilon=a_\varepsilon\beta_1+
       \sum_{j\geq2}\beta_j+\tau s\xi,
 \qquad a_\varepsilon=\varepsilon+2(1-\varepsilon)s^2.
\]
The first block of $\Omega_\varepsilon-\varepsilon\beta$ has
diagonal entries $2(1-\varepsilon)s^2$ and $1-\varepsilon$, and
determinant
\[
 \bigl[2(1-\varepsilon)^2-\tau^2\bigr]s^2\geq0.
\]
It is positive semidefinite; the remaining diagonal entries are
$1-\varepsilon$. Hence $\Omega_\varepsilon\geq\varepsilon\beta$,
with equality in the least eigenvalue at $s=0$.
Also $\Omega_\varepsilon\leq3\beta$, all fixed-order smooth norms
are uniformly bounded, and the smooth semipositive limit loses rank
at $s=0$. Since $d\ddc u_\varepsilon=0$, the exterior derivative
is the displayed fixed nonzero $dh_\tau$.

The lower-power expansion in the proof of \cref{thm:constant-mass-entire-determinant-profile} applies to $\Omega_\varepsilon=h_\tau+\ddc u_\varepsilon$; it gives $\ddc(\Omega_\varepsilon^j)=0$ for $j\leq n-1$, in particular pluriclosedness. Its volume is $V_\tau$, either by integrating the displayed determinant or by \cref{thm:mass-invariance-characterization}. That theorem also gives mass $V_\tau$ for every bounded $\Omega_\varepsilon$-plurisubharmonic potential, so both mass extrema equal $V_\tau$.
Theorem~\ref{thm:constant-mass-entire-determinant-profile} applied to $h_\tau$ gives equality of the
entire current and measure sets and their exact endpoint for every
$p\in L^n$. The balanced trace is fixed as well:
\[
 \int_X\Omega_\varepsilon\wedge\beta^{n-1}=n!.
\]
The off-diagonal term has zero trace and the mean of
$a_\varepsilon$ equals one. The global holomorphic volume form gives
$c_1=0$. The manifold itself is already K\"ahler, so these
nonclosed metrics show that the complete determinant profile and
positive bounded-potential mass extrema can stay fixed while the
least relative eigenvalue tends to zero. The coexistence problem
of \cref{qu:balanced-pluriclosed-program} remains open.
\end{proof}
